% Partie : logique du premier ordre
% Chapitre : complétude
% Section : ensembles maximalement cohérents

\documentclass[../../../include/open-logic-section]{subfiles}

\begin{document}

\olfileid[fr]{fol}{com}{mcs}
\olsection{Ensembles maximalement cohérents d'\usetoken{p}{sentence}}

\begin{defn}[Ensemble maximalement cohérent]
\ollabel{mcs}
Un ensemble~$\Gamma$ d'!!{sentence}s est
\emph{maximalement cohérent} si et seulement si :
\begin{enumerate}
\item $\Gamma$ est cohérent ;
\item tout ensemble $\Gamma'$ d'!!{sentence}s du même langage
tel que $\Gamma\subsetneq\Gamma'$ est incohérent.
\end{enumerate}
\end{defn}

\begin{explain}
Une définition équivalente est la suivante : un ensemble
$\Gamma$ d'énoncés est \emph{maximalement cohérent} si et
seulement si :
\begin{enumerate}
\item $\Gamma$ est cohérent ;
\item pour tout énoncé $!A$, si $\Gamma\cup\{!A\}$ est
cohérent, alors $!A\in\Gamma$.
\end{enumerate}
Autrement dit, on ne peut ajouter à un ensemble maximalement
cohérent~$\Gamma$ des !!{sentence}s qui n'y figurent pas déjà
sans rendre le nouvel ensemble incohérent.
\end{explain}

\begin{explain}
Les ensembles maximalement cohérents jouent un rôle important
dans la démonstration de complétude. Nous pouvons garantir que
tout ensemble cohérent d'!!{sentence}s~$\Gamma$ est inclus
dans un ensemble maximalement cohérent~$\Gamma^*$, et qu'un
tel ensemble contient, pour chaque !!{sentence}~$!A$, soit
$!A$, soit sa négation $\lnot !A$. Cela vaut notamment pour
les !!{sentence}s atomiques. À partir d'un ensemble
maximalement cohérent dans un langage convenablement enrichi
de !!{constant}s, nous pouvons donc construire !!a{structure}
dont l'interprétation des !!{predicate}s dépend des
!!{sentence}s atomiques appartenant à~$\Gamma^*$.
Nous pouvons ensuite montrer que cette !!{structure} rend
vrais tous les !!{sentence}s de~$\Gamma^*$, et donc aussi
ceux de~$\Gamma$. Pour établir ce dernier fait, il faut que
$\lnot !A\in\Gamma^*$ si et seulement si $!A\notin\Gamma^*$,
que $(!A\lor !B)\in\Gamma^*$ si et seulement si
$!A\in\Gamma^*$ ou $!B\in\Gamma^*$, et ainsi de suite.
\end{explain}

\begin{prop}\ollabel{prop:mcs}
Supposons que $\Gamma$ soit maximalement cohérent. Alors :
\begin{enumerate}
\item \ollabel{prop:mcs-prov-in}
Si $\Gamma\Proves !A$, alors $!A\in\Gamma$.

\item \ollabel{prop:mcs-either-or}
Pour tout énoncé $!A$, soit $!A\in\Gamma$, soit
$\lnot !A\in\Gamma$.

\tagitem{prvAnd}{\ollabel{prop:mcs-and}
$(!A\land !B)\in\Gamma$ si et seulement si
$!A\in\Gamma$ et $!B\in\Gamma$.}{}

\tagitem{prvOr}{\ollabel{prop:mcs-or}
$(!A\lor !B)\in\Gamma$ si et seulement si
$!A\in\Gamma$ ou $!B\in\Gamma$.}{}

\tagitem{prvIf}{\ollabel{prop:mcs-if}
$(!A\lif !B)\in\Gamma$ si et seulement si
$!A\notin\Gamma$ ou $!B\in\Gamma$.}{}
\end{enumerate}
\end{prop}

\begin{proof}
Dans toute la démonstration, supposons que $\Gamma$ soit
maximalement cohérent.
\begin{enumerate}
\item Si $\Gamma\Proves !A$, alors $!A\in\Gamma$.

Supposons que $\Gamma\Proves !A$ et, par l'absurde, que
$!A\notin\Gamma$. La maximalité de $\Gamma$ implique que
$\Gamma\cup\{!A\}$ est incohérent, donc que
$\Gamma\cup\{!A\}\Proves\lfalse$. D'après les
\tagrefs{prfSC/{fol:seq:prv:prop:provability-contr},prfND/{fol:ntd:prv:prop:provability-contr}},
$\Gamma$ est incohérent, contrairement à l'hypothèse.
Ainsi, $!A\notin\Gamma$ est impossible et $!A\in\Gamma$.

\item Pour tout énoncé $!A$, soit $!A\in\Gamma$, soit
$\lnot !A\in\Gamma$.

Supposons, par l'absurde, qu'un énoncé~$!A$ vérifie
$!A\notin\Gamma$ et $\lnot !A\notin\Gamma$.
Par maximalité, $\Gamma\cup\{!A\}$ et
$\Gamma\cup\{\lnot !A\}$ sont alors tous deux incohérents.
Ainsi, $\Gamma\cup\{!A\}\Proves\lfalse$ et
$\Gamma\cup\{\lnot !A\}\Proves\lfalse$. D'après les
\tagrefs{prfSC/{fol:seq:prv:prop:provability-exhaustive},prfND/{fol:ntd:prv:prop:provability-exhaustive}},
$\Gamma$ est incohérent, ce qui est contradictoire.
Il n'existe donc pas de tel !!{sentence}~$!A$ :
pour tout énoncé $!A$, on a $!A\in\Gamma$ ou
$\lnot !A\in\Gamma$.

\tagitem{defAnd}{}{%
\iftag{probAnd}{Exercice.}{%
$(!A\land !B)\in\Gamma$ si et seulement si
$!A\in\Gamma$ et $!B\in\Gamma$.

Pour le sens direct, supposons que $(!A\land !B)\in\Gamma$.
Alors $\Gamma\Proves !A\land !B$. D'après les
\tagrefs{prfSC/{fol:seq:ppr:prop:provability-land},prfND/{fol:ntd:ppr:prop:provability-land}}, point~(1),
$\Gamma\Proves !A$ et $\Gamma\Proves !B$.
La \olref{prop:mcs-prov-in} donne $!A\in\Gamma$ et
$!B\in\Gamma$, comme souhaité.

Réciproquement, supposons que $!A\in\Gamma$ et
$!B\in\Gamma$. Alors $\Gamma\Proves !A$ et
$\Gamma\Proves !B$. D'après les
\tagrefs{prfSC/{fol:seq:ppr:prop:provability-land},prfND/{fol:ntd:ppr:prop:provability-land}}, point~(2),
$\Gamma\Proves !A\land !B$. La \olref{prop:mcs-prov-in}
donne donc $(!A\land !B)\in\Gamma$.}}

\tagitem{defOr}{}{%
\iftag{probOr}{Exercice.}{%
$(!A\lor !B)\in\Gamma$ si et seulement si
$!A\in\Gamma$ ou $!B\in\Gamma$.

Pour la contraposée du sens direct, supposons que
$!A\notin\Gamma$ et $!B\notin\Gamma$. Nous voulons
montrer que $(!A\lor !B)\notin\Gamma$.
Par maximalité de $\Gamma$,
$\Gamma\cup\{!A\}\Proves\lfalse$ et
$\Gamma\cup\{!B\}\Proves\lfalse$. D'après les
\tagrefs{prfSC/{fol:seq:ppr:prop:provability-lor},prfND/{fol:ntd:ppr:prop:provability-lor}}, point~(1),
$\Gamma\cup\{(!A\lor !B)\}$ est incohérent.
Par conséquent, $(!A\lor !B)\notin\Gamma$, comme souhaité.

Réciproquement, supposons que $!A\in\Gamma$ ou
$!B\in\Gamma$. Alors $\Gamma\Proves !A$ ou
$\Gamma\Proves !B$. D'après les
\tagrefs{prfSC/{fol:seq:ppr:prop:provability-lor},prfND/{fol:ntd:ppr:prop:provability-lor}}, point~(2),
$\Gamma\Proves !A\lor !B$. La \olref{prop:mcs-prov-in}
donne donc $(!A\lor !B)\in\Gamma$, comme souhaité.}}

\tagitem{defIf}{}{%
\iftag{probIf}{Exercice.}{%
$(!A\lif !B)\in\Gamma$ si et seulement si
$!A\notin\Gamma$ ou $!B\in\Gamma$.

Pour le sens direct, supposons que $(!A\lif !B)\in\Gamma$.
Supposons, par l'absurde, que $!A\in\Gamma$ et
$!B\notin\Gamma$. Alors $\Gamma\Proves !A\lif !B$
et $\Gamma\Proves !A$. D'après les
\tagrefs{prfSC/{fol:seq:ppr:prop:provability-lif},prfND/{fol:ntd:ppr:prop:provability-lif}}, point~(1),
$\Gamma\Proves !B$. La \olref{prop:mcs-prov-in} donne
alors $!B\in\Gamma$, contrairement à l'hypothèse.

Pour la réciproque, considérons d'abord le cas
$!A\notin\Gamma$. Par la \olref{prop:mcs-either-or},
$\lnot !A\in\Gamma$, donc $\Gamma\Proves\lnot !A$.
D'après les
\tagrefs{prfSC/{fol:seq:ppr:prop:provability-lif},prfND/{fol:ntd:ppr:prop:provability-lif}}, point~(2),
$\Gamma\Proves !A\lif !B$. La \olref{prop:mcs-prov-in}
donne de nouveau $(!A\lif !B)\in\Gamma$, comme souhaité.

Considérons maintenant le cas $!B\in\Gamma$.
Alors $\Gamma\Proves !B$ ; d'après les
\tagrefs{prfSC/{fol:seq:ppr:prop:provability-lif},prfND/{fol:ntd:ppr:prop:provability-lif}}, point~(2),
$\Gamma\Proves !A\lif !B$. La \olref{prop:mcs-prov-in}
donne donc $(!A\lif !B)\in\Gamma$.}}
\end{enumerate}
\end{proof}

\begin{prob}
Achever la démonstration de la \olref[fol][com][mcs]{prop:mcs}.
\end{prob}

\end{document}
